\hypertarget{appendix-49933872bb51}{}\label{referencestart:appendix-49933872bb51}
\section{Norms lengths and induced distances}\label{appendix-49933872bb51-norms-lengths-and-induced-distances}

\textbf{Curriculum depth:} Introduced · \textbf{First seen:} Phase6 · \textbf{Latest development:} Phase6

\begin{allmightroledefinition}

\subsection{Definition and mechanism}\label{appendix-49933872bb51-definition-and-mechanism}

A norm is nonnegative, definite, absolutely homogeneous and satisfies the triangle inequality. It induces distance by taking the norm of a point difference. Coordinate-defined first, second and infinity norms answer different geometric questions.

\end{allmightroledefinition}

For \((3,-4)\), the three norms are \(7\), \(5\) and \(4\). Squared length scales quadratically and is not itself a norm or induced metric.

\subsection{Prerequisites}\label{appendix-49933872bb51-prerequisites}

\begin{itemize}
\tightlist
\item
  \hyperref[appendix-ce0341f7ea0d]{Basis coordinates and dimension} --- An ordered basis fixes coordinate meaning and the space whose geometry is measured.
\end{itemize}

\subsection{Used by}\label{appendix-49933872bb51-used-by}

\begin{itemize}
\item
  \hyperref[appendix-f744e41b95b4]{Unit directions and safe normalization} --- Builds on this mechanism in the connected geometry story.
\item
  \href{../06-GPU-Architecture-and-Real-Time-Rendering/main.pdf\#nameddest=appendix-8bc4e8fc547c}{Dot products and orthogonality (GPU Architecture and Real-Time Rendering, Reference notes: Dot products and orthogonality)} --- Builds on this mechanism in the connected geometry story.
\item
  \href{../05-Spatial-Worlds-Engine-State-and-Simulation/main.pdf\#nameddest=appendix-f8657bac66e5}{Matrix norms and numerical error (Spatial Worlds, Engine State, and Simulation, Reference notes: Matrix norms and numerical error)} --- Extends this foundation into linear-systems and decomposition analysis.
\item
  \href{../04-Light-Transport-and-Image-Synthesis/main.pdf\#nameddest=appendix-a9ba01fa72a6}{Collision barriers and continuous detection (Light Transport and Image Synthesis, Reference notes: Collision barriers and continuous detection)} \#\# Recall and next depth
\end{itemize}

Draw the three unit balls; explain which comparisons can safely use squared length. Full production proofs, target-specific integration and learner mastery require further evidence.

\subsection{Source and navigation}\label{appendix-49933872bb51-source-and-navigation}

\href{../02-Mathematics-for-Visual-Computing/main.pdf\#nameddest=01-foundations-i-ch017-a-ruler-needs-more-than-coordinate-components}{Teaching narrative (Mathematics for Visual Computing, Ch. 1)} · \hyperref[appendix-fd090d982a15-phase-6---parts-501-600]{Coverage and source gap} · \hyperref[appendix-a103e7846d21-phase-6---source-analysis-and-verification]{Evidence and limits} · \hyperref[appendix-4ad47bdd269a]{Dependency Map}
